\section{Thirteen beta values and a rational Gram block}\label{sec:gram}
This construction has a different advantage: its polynomial degree stays
bounded while its matrix size grows. It involves precisely
$\beta(4),\beta(6),\ldots,\beta(28)$. The question is whether the growing
matrix supplies useful decay after its rational part has been removed.
The one-label antecedent is the ordinary-moment Schur construction in
\cite[Section~8]{firstsupp}; we add seven logarithmic labels.

\subsection{The seven boundary functions}
On $0<x<1$, put $L=-\log x$, $t=x^2$, and
$\dd\mu=L^3\dd x/(1+x^2)$. For $r\geq1$, take the Gram matrix of
\[
 t^jL^{2a},\qquad 0\leq a\leq6,\quad0\leq j<r.
\]
These functions are independent, as is seen by substituting $x=e^{-y}$
and using independence of exponential polynomials. The matrix is positive
definite. For $s=4+2a+2b$, its entries are
\[
 (s-1)!(-1)^{j+k}\left(\beta(s)-
          \sum_{v=0}^{j+k-1}\frac{(-1)^v}{(2v+1)^s}\right).
\]
The formula follows by a finite geometric identity and integration
of $x^m(-\log x)^{s-1}$.

For each logarithmic label replace
$1,t,\ldots,t^{r-1}$ by $1,(1+t),(1+t)t,\ldots,(1+t)t^{r-2}$.
This change has determinant one. Grouping the seven boundary functions
$L^{2a}$ first gives
\begin{equation}\label{eq:seven-schur}
 \begin{pmatrix}A&C\\C^T&B\end{pmatrix},\qquad
 S_r=A-CB^{-1}C^T,\qquad \det G_r=\det B\det S_r.
\end{equation}
Write $X_k$ for $\beta(4+2k)$, $0\leq k\leq12$. The blocks are
\begin{align*}
 A_{ab}&=(3+2a+2b)!X_{a+b},\\
 C_{a,(b,j)}&=\frac{(3+2a+2b)!}{(2j+1)^{4+2a+2b}},\\
 B_{(a,i),(b,j)}&=(3+2a+2b)!\left(
 (2i+2j+1)^{-4-2a-2b}+(2i+2j+3)^{-4-2a-2b}\right),
\end{align*}
where $0\leq i,j\leq r-2$. For $r=1$ the bulk is empty.
$B$ is rational positive definite, so the primitive polynomial associated
with $\det G_r$ is that associated with
$P_r(X)=\det S_r(X)$. The rational factor $\det B$ supplies no separate
gain. This exact reduction alone does not determine the residual decay.

The one-label case at $r=2$ makes the normalization visible:
$A=6\bet$, $C=6$, and $B=164/27$, so
\[
 S=6\bet-\frac{243}{41}
   =\frac3{41}(82\bet-81).
\]
The primitive polynomial is $82X-81$. The factor from the rational
bulk cannot make its primitive value smaller.

\subsection{The residual decays on a linear scale}
Let $\rho=3+\sqrt8$ and let $T_j$ denote the Chebyshev polynomial.
Define
\[
 H_0=\prod_{a=0}^6(3+4a)!,\quad
 V_0=\pi^{84}\left(\prod_{j=1}^6(2j)!\right)^2,
 \quad M_0=\sqrt{1+(13\pi/2)^2},\quad
 \Gamma=28+28\sqrt{26M_0}<700.
\]
\begin{theorem}\label{thm:residual}
At the actual beta values, for every $r\geq1$,
\begin{equation}\label{eq:residualbounds}
 \frac{V_0e^{-\Gamma r}}{(52r^2)^7}
 \leq P_r(\beta)
 \leq\frac{H_0}{T_{r-1}(3)^{14}}
 \leq2^{14}H_0\rho^{-14(r-1)}.
\end{equation}
In particular $\log P_r(\beta)=-\Theta(r)$ and
$r^{-2}\log P_r(\beta)\longrightarrow0$.
\end{theorem}
\begin{proof}
The Schur complement expresses projection onto the orthogonal complement
of the bulk. For any real $v=(v_a)$,
\[
 v^TS_rv=\min_{p_a(-1)=v_a,\ \deg p_a<r}
       \int_0^1\left|\sum_{a=0}^6p_a(x^2)L^{2a}\right|^2\dd\mu.
\]
The constraint says exactly that $p_a(t)=v_a-(1+t)Q_a(t)$ for a bulk
polynomial $Q_a$. Choose
$p_a(t)=v_aT_{r-1}(2t-1)/T_{r-1}(-3)$.
Its common factor is at most $1/T_{r-1}(3)$ on $[0,1]$. Hence
$0<S_r\leq A/T_{r-1}(3)^2$ in matrix order. Take determinants,
apply Hadamard to $A$, and use $\beta(s)<1$ to obtain the upper bound.
The formula $T_j(3)=(\rho^j+\rho^{-j})/2$ gives the final inequality.

For the lower bound, write $y=-\log x$. Every residual lies in
\[
 \mathcal E_r=\Span\{y^he^{-2jy}:0\leq j<r,\ 0\leq h\leq12\},
 \quad
 \|F\|_\mu^2=\int_0^\infty
       |F(y)|^2\frac{y^3e^{-y}}{1+e^{-2y}}\dd y.
\]
At the seven points $z_l=i\pi(l+1/2)$, the bulk factor $1+e^{-2z_l}$
vanishes. The residual $R_a$ of the boundary function $y^{2a}$
therefore satisfies $R_a(z_l)=z_l^{2a}$, and
\begin{equation}\label{eq:boundaryV}
 |\det(R_a(z_l))_{l,a=0}^6|^2=V_0.
\end{equation}
This is a Vandermonde identity: $\prod_{i<j}(j-i)(j+i+1)=
\prod_{j=1}^6(2j)!$.

The repeated-exponent evaluation estimate proved in
Appendix~\ref{app:evaluation} gives, for $F\in\mathcal E_r$,
\[
 |F(z_l)|^2\leq52r^2e^{\Gamma r/7}\|F\|_\mu^2.
\]
Orthonormalize the seven residuals. Each row of their evaluation matrix
has squared length at most this evaluation-norm bound. Hadamard bounds
its squared determinant by the seventh power. Undoing the basis change
and using \eqref{eq:boundaryV} gives the lower bound in
\eqref{eq:residualbounds}.
\end{proof}

\subsection{The arithmetic question left by this reduction}
Let $q_r>0$ make $q_rP_r$ primitive integral. Its total degree is seven.
The coefficient of $X_6^7$ is $-(15!)^7$: only the reversing permutation
$a\mapsto6-a$ contributes. Thus $q_r\geq1/(15!)^7$.
Under joint rationality of the thirteen values, a fixed common
denominator raised to the seventh power would make every positive
$q_rP_r(\beta)$ an integer. It is therefore sufficient to prove
$q_rP_r(\beta)\to0$ on an unbounded subsequence.

There are useful exact bounds in both directions. Let $\ell(m)$ be
the lcm of the positive odd integers at most $m$. For $r\geq2$, put
\[
 H_{aa}=\ell(2r-3)^{2+2a},\qquad
 E_{(a,i),(a,i)}=\ell(2i+2r-1)^{2+2a}.
\]
The matrices $M=EBE$, $N=HCE$ and $HA(X)H$ are integral: split each
denominator exponent $4+2a+2b$ between its two indices. The block
determinant
\[
 \det\begin{pmatrix}HAH&N\\N^T&M\end{pmatrix}
       =U_rP_r(X),\qquad U_r=(\det H)^2\det M\in\Z_{>0}
\]
therefore gives $q_r\leq U_r$. Hadamard and
$\sum_{a=0}^6(2+2a)=56$ imply
\[
 U_r\leq C_0^{r-1}
 \left(\ell(2r-3)\prod_{i=0}^{r-2}\ell(2i+2r-1)\right)^{112},
 \quad C_0=\prod_{a=0}^6 2(3+4a)!.
\]
The prime number theorem gives $\log q_r\leq336r^2+o(r^2)$.
This upper bound is too large for the linear-scale sufficient condition,
but says nothing about whether the true cost is quadratic.

For a necessary condition, put $R_r=CB^{-1}C^T$ and
$F_7=\prod_{a=1}^6(3+4a)!$. The degree-six coefficient
\begin{equation}\label{eq:R00}
 [X_2X_4X_6X_8X_{10}X_{12}]P_r=-F_7(R_r)_{00}
\end{equation}
isolates one rational entry. Give $X_k$ weight $k$. The cofactor of
entry $(i,j)$ of $A$ has weight $42-i-j$; only $(0,0)$ can contribute
to the displayed weight $42$. In that cofactor the successively
smallest variables force the diagonal, proving \eqref{eq:R00}.
If $b_r$ is the reduced denominator of $(R_r)_{00}$ and
$c_7=(15!)^7$, then integrality of both coefficients gives
\[
 q_r\geq\frac{b_r}{c_7\gcd(b_r,F_7)}.
\]
A superlinear lower bound for $\log b_r$ would rule out successful
shrinking on that sequence. None has been proved here.

Along any chosen unbounded sequence of indices,
Theorem~\ref{thm:residual} gives the sufficient condition
$\limsup\log q_r/r<14\log\rho=24.67846\ldots$, whereas convergence
of the primitive values to zero requires
$\limsup\log q_r/r\leq\Gamma$. Neither arithmetic bound decides
this comparison.

\begin{center}
\begin{tabular}{rrrr}
\toprule
$r$&$\log q_r$&$\log P_r(\beta)$&$\log(q_rP_r(\beta))$\\
\midrule
1&$-152.01$&194.93&42.92\\
2&14.04&$-3.33$&10.71\\
3&655.17&$-119.25$&535.92\\
4&1586.31&$-254.99$&1331.32\\
5&3103.65&$-402.44$&2701.21\\
6&4842.52&$-545.27$&4297.25\\
\bottomrule
\end{tabular}
\end{center}
These certified finite values all exceed one after primitive
normalization. They illustrate the arithmetic issue, but do not prove
growth along every unbounded subsequence. The general conclusion is
the residual scale theorem, with primitive arithmetic unresolved.
